\subsection{Erreur Quadratique Moyenne}
L'Erreur Quadratique Moyenne (EQM, ou \textit{Mean Squared Error} : MSE en anglais) 
($MSE = \frac{1}{n} \sum(y - \hat{y})^2$) 
mesure la moyenne des carrés des écarts entre les valeurs prédites et les valeurs 
réelles. En élevant au carré les erreurs de prédiction, la MSE pénalise les 
erreurs importantes. 


\subsection{Erreur Absolue Moyenne (MAE) }

L'Erreur Absolue Moyenne ($MAE= \frac{1}{n} \sum|y - \hat{y}|$ ) est une métrique 
qui mesure, en moyenne, la magnitude des erreurs entre les valeurs prédites par 
le modèle et les valeurs 
réelles observées, sans tenir compte de la direction de l'erreur. 
La MAE est appréciée pour sa simplicité et son interprétation intuitive. Elle donne 
une estimation directe de l'erreur moyenne sur l'ensemble des prédictions. 

\subsection{Mean\_poisson\_deviance}
Elle mesure l'erreur entre les valeurs réelles $y$ et les prédictions 
$\hat{y}$ pour des données de comptage (0,1,2,\dots). Plus la valeur est petite, 
meilleur est le modèle.
Pour une observation :
\[
D(y, \hat{y}) = 2 \left( y \log\left(\frac{y}{\hat{y}}\right) - (y - \hat{y}) \right)
\]
Pour $n$ observations :
\[
\text{mean\_poisson\_deviance} = \frac{1}{n} \sum_{i=1}^{n} 2 \left( y_i \log\left(\frac{y_i}{\hat{y}_i}\right) - (y_i - \hat{y}_i) \right)
\]


